\section{Polylogarithmic endpoint identities and \texorpdfstring{$z$}{z}-primitives}\label{sec:polylog}

\subsection{Regularized hypergeometric notation}
For $|z|<1$ define
\begin{equation}\label{eq:F}
 F(u,z)=\sum_{n\ge0}r_n(u)z^n
 =\Gamma(a)\Gamma(b)\,{}_2\mathbf F_1(a,b;a+b+u;z),
\end{equation}
where ${}_2\mathbf F_1(a,b;d;z)={}_2F_1(a,b;d;z)/\Gamma(d)$
is the \emph{regularized} function, defined by its coefficient series.
This convention is indispensable when $d$ is a nonpositive integer:
separate floating-point evaluation of an ordinary hypergeometric pole and
a reciprocal-gamma zero is not the definition.

The Lerch transcendent satisfies
\begin{equation}\label{eq:Lerch}
 \Phi(z,s,c)=\sum_{n\ge0}\frac{z^n}{(n+c)^s},\qquad
 z\Phi(z,s,1)=\Li_s(z),
\end{equation}
initially on $|z|<1$ \cite{dlmf-lerch}.

\begin{theorem}[Polylogarithmic endpoint subtraction, with all jets]\label{thm:endpoint}
Let $K\ge1$.  The combined function
\begin{equation}\label{eq:FK}
 F_K(u,z;c)=F(u,z)-\sum_{j=0}^{K-1}
                         A_j(u;c)\Phi(z,1+u+j,c)
\end{equation}
has the absolutely convergent coefficient representation
\[
 F_K(u,z;c)=\sum_{n\ge0}\left\{r_n(u)
       -\sum_{j=0}^{K-1}\frac{A_j(u;c)}{(n+c)^{1+u+j}}\right\}z^n
\]
for every $|z|\le1$ and $\Rea u>-K$.
It is continuous in $z$ on the closed disk and holomorphic in $u$.
The convergence and every fixed-order $u$ derivative are uniform on the
closed disk and compact subsets of that half-plane.  In particular,
\begin{equation}\label{eq:endpoint}
 \lim_{z\to1,\ |z|\le1}[t^m]F_K(-N+t,z;c)
 =S_{N,m}^{(K)}(c),\qquad N<K.
\end{equation}
For $c=1$ this is an endpoint identity with finitely many ordinary
polylogarithms and their spectral derivatives:
\begin{equation}\label{eq:polyendpoint}
 zF_K(u,z;1)=zF(u,z)-\sum_{j=0}^{K-1}
                                    A_j(u;1)\Li_{1+u+j}(z).
\end{equation}
\end{theorem}
\begin{proof}
The coefficient bracket is the one in Theorem~\ref{thm:subtraction}.
The same summable majorant works for every $|z|\le1$, since $|z^n|\le1$.
The Weierstrass test, also after a fixed number of spectral derivatives,
proves uniform convergence, continuity, and interchange of endpoint limits
with coefficient extraction.  Equations \eqref{eq:FK} and
\eqref{eq:polyendpoint} first hold inside the disk by absolutely convergent
power series, and then represent the claimed continuous extension of the
\emph{combined} expression.
\end{proof}

\begin{remark}[The separate pieces can still be singular]
The theorem does not assert that each Lerch or polylogarithm term has a
finite endpoint separately.  Nor does it rely on identifying two
independently regularized boundary values.  The coefficient differences
already form an ordinary uniformly convergent series, which fixes the
limit uniquely.  The spectral-derivative formula is therefore meaningful
at negative integer polylogarithm orders and at order one alike.
\end{remark}

\subsection{Explicit endpoint examples}
In the half-parameter case, write
$\Fcal(z)=\sum_{n\ge0}w_nz^n=\pi{}_2F_1(\tfrac12,\tfrac12;1;z)$
and $\theta=z\partial_z$.  For $N\ge0$, the resonant series is
$F(-N,z)=\theta(\theta-1)\cdots(\theta-N+1)\Fcal(z)$.
For instance,
\begin{align}
 \lim_{z\to1^-}\{z\Fcal(z)-\Li_1(z)\}
 &=4\log2,\label{eq:z0}\\
 \lim_{z\to1^-}\{z\theta\Fcal(z)-\Li_0(z)+\tfrac14\Li_1(z)\}
 &=-\log2-\tfrac14,\label{eq:z1}\\
 \lim_{z\to1^-}\{z\theta(\theta-1)\Fcal(z)-\Li_{-1}(z)
       +\tfrac94\Li_0(z)-\tfrac9{32}\Li_1(z)\}
 &=\tfrac98\log2+\tfrac{27}{64}.\label{eq:z2}
\end{align}
These are ordinary limits of the displayed expressions.  They are
instances of the all-$N$ formula, not numerical extrapolations.

The first order-derivative example is
\begin{align}
 \lim_{z\to1^-}\left\{
 z\left.\partial_uF(u,z)\right|_{u=0}
 -\left.\partial_s\Li_s(z)\right|_{s=1}\right\}
 =\frac{(\gamma+4\log2)^2-5\zeta(2)}2+\gamma_1.
 \label{eq:zstieltjes}
\end{align}
All higher orders and all negative resonances follow by the explicit finite
formula \eqref{eq:alljets}.  This is the direct connection between the
polylogarithmic order calculus and the Stieltjes layer in this construction.

\subsection{Exact Euler derivatives and normalized antiderivatives}
The operator $\theta$ multiplies a coefficient of $z^n$ by $n$.
After $q$ applications, the endpoint majorant becomes
$O(n^{q-1-\Rea u-K}\log^m n)$.  Thus the same endpoint statement for
$\theta^qF_K$ is valid on the smaller half-plane
\begin{equation}\label{eq:thetadomain}
 \Rea u>q-K.
\end{equation}
One cannot differentiate the boundary formula arbitrarily often while
keeping the original number of counterterms.

For $c>0$, define the zero-basepoint normalized primitive
\begin{equation}\label{eq:Ic}
 (\mathcal I_cf)(z)=z^{-c}\int_0^z t^{c-1}f(t)\,\dd t.
\end{equation}
It is first interpreted for real $0<z<1$, and then by its analytic power
series.  Coefficientwise, $\mathcal I_c(z^n)=z^n/(n+c)$; hence
\begin{equation}\label{eq:IcLerch}
 \mathcal I_c^{\ell}\Phi(z,s,c)=\Phi(z,s+\ell,c),\qquad \ell\ge0.
\end{equation}
For the gamma-ratio term this gives the exact lift
\begin{align}
 \mathcal I_c^{\ell}F(u,z)
 ={}&\frac{\Gamma(a)\Gamma(b)}{\Gamma(a+b+u)c^{\ell}}
 {}_{\ell+2}F_{\ell+1}
 \left(\begin{matrix}a,b,\underbrace{c,\ldots,c}_{\ell}\\
 a+b+u,\underbrace{c+1,\ldots,c+1}_{\ell}\end{matrix};z\right).
 \label{eq:hyperlift}
\end{align}
As in \eqref{eq:F}, reciprocal-gamma regularization is understood at an
exceptional lower parameter.  The proof is the elementary identity
$(c)_n/(c+1)_n=c/(n+c)$ in every coefficient.

Applying \eqref{eq:IcLerch}--\eqref{eq:hyperlift} to \eqref{eq:FK} provides
exact antiderivatives of all the endpoint-subtracted families.
Spectral differentiation commutes with these primitives inside the disk,
and with their endpoint limits whenever the corresponding coefficient
majorant is summable.  This is an integration ladder, not a claim that
arbitrary higher hypergeometric endpoint values reduce to gamma values.
Likewise, the generic function $F(u,z)$ is not asserted to be an ordinary
polylogarithm.  The proved polylogarithmic statement is the finite
subtraction and the exact evaluation of its combined endpoint.
